\documentclass[11pt]{article}

% Mathematics
\usepackage{amsmath}
\usepackage{amssymb}
\usepackage{amsthm}
\usepackage{mathtools}

% Fonts / symbols
\usepackage{mathrsfs}

% Figures
\usepackage{graphicx}

% Tables
\usepackage{booktabs}

% Lists
\usepackage{enumitem}

% TikZ / diagrams
\usepackage{tikz}
\usetikzlibrary{
    arrows.meta,
    positioning,
    shapes,
    calc
}

\newtheorem{theorem}{Theorem}
\newtheorem{lemma}[theorem]{Lemma}
\newtheorem{proposition}[theorem]{Proposition}

\theoremstyle{definition}
\newtheorem{definition}[theorem]{Definition}
\newtheorem{example}[theorem]{Example}

\theoremstyle{remark}
\newtheorem{remark}[theorem]{Remark}

% URLs and clickable references
\usepackage{hyperref}

% Better references: \cref{...}
\usepackage{cleveref}

% Bibliography
\usepackage[
    backend=biber,
    style=alphabetic
]{biblatex}
\addbibresource{references.bib}

\begin{document}


\begin{definition}
    Let $Var$ be the set of propositional variables. Then $\phi$ is defined inductively as follows: 
    \[
       \phi ::=  p \mid \phi \lor \phi \mid X\phi \mid  \phi U \psi \mid \langle a \rangle \phi 
    \]
\end{definition}

\begin{definition}
    Let  $\phi $ be a formula and $p \in Var$. Then the Closure of $\phi$ is inductively defined as

\[
    \begin{aligned}
        Cl(p) &= Cl(\neg p) = \{p,\neg p\} \\
        Cl(\neg \phi) &= Cl(\phi) \\
        Cl(\phi \lor \psi) &= \{\phi \lor \psi, \neg(\phi \lor \psi)\} \cup  Cl(\phi) \cup Cl(\psi) \\
        Cl(X \phi) &= \{X \phi, \neg X \phi\} \cup Cl(\phi) \\
        Cl(\phi U \psi) &= \{\phi U \psi, \neg(\phi U \psi)\} \cup Cl(\phi) \cup Cl(\psi) \\
        Cl(\langle a \rangle \phi) &= \{\langle a \rangle \phi, \neg \langle a \rangle \phi\} \cup Cl(\phi)
    \end{aligned}
\]
\end{definition}

\begin{definition}
An atom $A$ is a subset of $Cl(\phi)$ such that
\[
\begin{aligned}
&\text{if } \psi \in Cl(\phi),
&&\text{ then } \psi \in A \text{ iff } \neg \psi \notin A, \\
&\text{if } \phi_1 \lor \phi_2 \in Cl(\phi), 
&&\text{ then } \phi_1 \lor \phi_2 \in A\text{ iff } \phi_1 \in A \text{ or } \phi_2 \in A \\
&\text{if }  \phi_1 U \phi_2 \in Cl(\phi),
&&\text{then  if } \phi_1 U \phi_2 \in A \mbox{ then } \psi_1 \in A \mbox{ or } \psi_2 \in A \\   
&\text{if }  \phi_1 U \phi_2 \in Cl(\phi),
&&\text{then  if } \phi_2 \in A \mbox{ then } \phi_1 U \phi_2 \in A
\end{aligned}
\]
\end{definition}

\begin{definition}
    We say two atoms $A,B$ are in relation $R_x (A R_x B)$ iff

\[
\begin{aligned}
&\text{if } X\psi \in Cl(\phi),
&&\text{ then } X\psi \in A \text{ iff } \psi \in B, \\[2mm]
&\text{if } \psi_1 U \psi_2 \in Cl(\phi),
&&\text{ then if } \psi_1 U \psi_2, \neg\psi_2 \in A,
   \text{ then } \psi_1 U \psi_2 \in B, \\
&&&\text{ and if } \psi_1 U \psi_2 \in B, \psi_1 \in A,
   \text{ then } \psi_1 U \psi_2 \in A.
\end{aligned}
\]


\end{definition}

\end{document}